%!TEX root = ../../note.tex
% Source: ../note-course/chapters/u2/s3.tex
\subsection{\texorpdfstring{$\R$}{R} as an ordered field}
The real numbers form a \emph{complete ordered field}. We build this up in
three stages: the field axioms, which govern $+$ and $\cdot$; the order
axioms, which let us compare elements; and the completeness axiom, which
distinguishes $\R$ from $\Q$ and is the source of all limiting behaviour in
analysis.

\subsubsection*{The field axioms}
\begin{defi}[Field]
  A \term{field} is a set $F$ with two binary operations
  $+ : F \times F \to F$ and $\cdot : F \times F \to F$ such that for all
  $a, b, c \in F$:
  \begin{enumerate}[label=\textbf{F\arabic*.}, leftmargin=2.5em]
    \item (Associativity of addition) $a + (b + c) = (a + b) + c$.
    \item (Commutativity of addition) $a + b = b + a$.
    \item (Additive identity) There is $0 \in F$ with $a + 0 = a$.
    \item (Additive inverses) For every $a \in F$ there is $-a \in F$ with $a + (-a) = 0$.
    \item (Associativity of multiplication) $a \cdot (b \cdot c) = (a \cdot b) \cdot c$.
    \item (Commutativity of multiplication) $a \cdot b = b \cdot a$.
    \item (Multiplicative identity) There is $1 \in F$, $1 \neq 0$, with $a \cdot 1 = a$.
    \item (Multiplicative inverses) For every $a \in F$ with $a \neq 0$ there is $a^{-1} \in F$ with $a \cdot a^{-1} = 1$.
    \item (Distributive law) $a \cdot (b + c) = a \cdot b + a \cdot c$.
  \end{enumerate}
\end{defi}

\begin{eg}
  $\Q$ also satisfies F1--F9. So being a field says nothing about size,
  order or completeness, and does not distinguish $\R$ from $\Q$.
\end{eg}

Once the axioms are fixed, even the most familiar facts become theorems.
\begin{thm}\leavevmode
  \begin{enumerate}[label=(\alph*)]
    \item If $z, a \in \R$ and $z + a = a$, then $z = 0$.
    \item If $u, b \in \R$, $b \neq 0$ and $u \cdot b = b$, then $u = 1$.
    \item If $a \in \R$, then $a \cdot 0 = 0$.
  \end{enumerate}
\end{thm}

\begin{proof}\leavevmode
  \begin{enumerate}[label=(\alph*)]
    \item $z = z + 0 = z + (a - a) = (z + a) + (-a) = a + (-a) = 0$.
    \item Multiplying $u \cdot b = b$ by $b^{-1}$,
      $u = u \cdot 1 = u \cdot (b \cdot b^{-1}) = (u \cdot b) \cdot b^{-1} = b \cdot b^{-1} = 1$.
    \item $a \cdot 0 = a \cdot (0 + 0) = a \cdot 0 + a \cdot 0$. Adding
      $-(a \cdot 0)$ to both sides and using associativity gives
      $0 = a \cdot 0$.\qedhere
  \end{enumerate}
\end{proof}

\begin{remark}
  Every ``obvious'' manipulation is an axiom in disguise: moving a term
  across ``$=$'' uses additive inverses and associativity, cancelling a
  nonzero factor uses multiplicative inverses, and expanding a product uses
  the distributive law.
\end{remark}

\subsubsection*{The order axioms}
\begin{defi}[Order properties of $\R$]
  There is a nonempty subset $P \subseteq \R$, the set of \term{positive real
  numbers}, such that:
  \begin{enumerate}[label=\textbf{O\arabic*.}, leftmargin=2.5em]
    \item (Trichotomy) For every $a \in \R$, exactly one of $a = 0$,
      $a \in P$, $-a \in P$ holds.
    \item (Closure under addition) If $a, b \in P$, then $a + b \in P$.
    \item (Closure under multiplication) If $a, b \in P$, then $a \cdot b \in P$.
  \end{enumerate}
\end{defi}

\begin{defi}[Order relation]
  For $a, b \in \R$, we write $a > b$ (or $b < a$) if $a - b \in P$, and
  $a \geq b$ (or $b \leq a$) if $a - b \in P \cup \{0\}$.
\end{defi}
So to show $a > b$, one shows $a - b \in P$; every manipulation of
inequalities traces back to this and to O1--O3.

The positive rationals $P \cap \Q$ also satisfy O1--O3, so an ordered field
still need not be $\R$. What $\Q$ lacks is completeness: the set
$\{q \in \Q : q^2 < 2\}$ is bounded above in $\Q$ but has no least upper
bound in $\Q$. This is why completeness is a separate axiom.

\begin{axiom}[Completeness axiom]
  If $A \subseteq \R$ is nonempty and bounded above, then $A$ has a least
  upper bound in $\R$.
\end{axiom}
Upper bounds, suprema and the consequences of this axiom are developed in
the next two subsections.

\begin{eg}
  If $a \in \R$ and $a \neq 0$, then $a^2 > 0$.
\end{eg}

\begin{proof}
  Since $a \neq 0$, trichotomy leaves two cases. If $a \in P$, then
  $a^2 = a \cdot a \in P$ by O3. Otherwise $-a \in P$, so
  $(-a)(-a) \in P$ by O3, and it remains to check $(-a)(-a) = a^2$.
  Multiplying $a + (-a) = 0$ by $-a$,
  \begin{align*}
    (-a)\bigl(a + (-a)\bigr) = (-a) \cdot 0 = 0
    &\implies (-a)a + (-a)(-a) = 0\\
    &\implies (-a)(-a) = -\bigl((-a)a\bigr) = -(-a^2) = a^2.
  \end{align*}
  In either case $a^2 \in P$, i.e.\ $a^2 > 0$.
\end{proof}
This is the typical use of trichotomy: split into the cases $a = 0$,
$a \in P$, $-a \in P$, and use O2 and O3 to keep track of positivity.

\subsubsection*{Absolute value and the triangle inequality}
\begin{defi}[Absolute value]
  For $a \in \R$, the \term{absolute value} of $a$ is
  \[
    \abs{a} =
    \begin{cases}
      a & a \geq 0,\\
      -a & a < 0.
    \end{cases}
  \]
\end{defi}
Geometrically, $\abs{a}$ is the distance from $a$ to $0$, and $\abs{a - b}$
is the distance between $a$ and $b$.

\begin{nthm}\label{thm:abs}
  For all $a, b \in \R$ and $c > 0$:
  \begin{enumerate}[label=(\alph*)]
    \item $\abs{ab} = \abs{a}\abs{b}$;
    \item $\abs{a}^2 = a^2$;
    \item $\abs{a} \leq c \iff -c \leq a \leq c$;
    \item $-\abs{a} \leq a \leq \abs{a}$.
  \end{enumerate}
\end{nthm}
Every part is proved the same way: split on the sign of $a$ (and of $b$) and
unwind the definition.

\begin{proof}
  The definition gives $\abs{-a} = \abs{a}$ and $-\abs{a} \leq a \leq \abs{a}$,
  which is (d); checking the cases $a \geq 0$ and $a < 0$ gives (b).

  For (a), if $a, b \geq 0$, then $ab \geq 0$ and $\abs{ab} = ab = \abs{a}\abs{b}$.
  If $a \geq 0 > b$, then $ab \leq 0$ and $\abs{ab} = -ab = a(-b) = \abs{a}\abs{b}$;
  the case $b \geq 0 > a$ is symmetric. If $a, b < 0$, then $ab > 0$ and
  $\abs{ab} = ab = (-a)(-b) = \abs{a}\abs{b}$.

  For (c), suppose $\abs{a} \leq c$. If $a \geq 0$, then $a = \abs{a} \leq c$
  and $-c < 0 \leq a$. If $a < 0$, then $-a = \abs{a} \leq c$, i.e.\
  $-c \leq a$, and $a < 0 < c$. Conversely, suppose $-c \leq a \leq c$. If
  $a \geq 0$, then $\abs{a} = a \leq c$. If $a < 0$, multiplying $-c \leq a$
  by $-1$ gives $\abs{a} = -a \leq c$.
\end{proof}

\begin{nprop}[Triangle inequality]\label{prop:triangle}
  If $a, b \in \R$, then $\abs{a + b} \leq \abs{a} + \abs{b}$.
\end{nprop}
The idea is to add the two-sided bounds of Theorem~\ref{thm:abs}(d) for $a$
and $b$, and then use part (c) to turn the result back into one inequality.

\begin{proof}
  By Theorem~\ref{thm:abs}(d), $-\abs{a} \leq a \leq \abs{a}$ and
  $-\abs{b} \leq b \leq \abs{b}$. Adding,
  \[
    -\bigl(\abs{a} + \abs{b}\bigr) \leq a + b \leq \abs{a} + \abs{b}.
  \]
  If $\abs{a} + \abs{b} = 0$, this forces $a + b = 0$ and the claim is
  immediate. Otherwise $\abs{a} + \abs{b} > 0$, and
  Theorem~\ref{thm:abs}(c) with $c = \abs{a} + \abs{b}$ gives
  $\abs{a + b} \leq \abs{a} + \abs{b}$.
\end{proof}
In terms of distance: going from $0$ to $a + b$ directly is never longer than
going via $a$.

\begin{cor}
  For all $a, b \in \R$:
  \begin{enumerate}
    \item $\bigl|\abs{a} - \abs{b}\bigr| \leq \abs{a - b}$;
    \item $\abs{a - b} \leq \abs{a} + \abs{b}$.
  \end{enumerate}
\end{cor}

\begin{proof}\leavevmode
  \begin{enumerate}
    \item Write $a = (a - b) + b$. The triangle inequality gives
      $\abs{a} \leq \abs{a - b} + \abs{b}$, i.e.\
      $\abs{a} - \abs{b} \leq \abs{a - b}$. Symmetrically, $b = (b - a) + a$
      gives $\abs{b} - \abs{a} \leq \abs{b - a} = \abs{a - b}$. Together,
      \[
        -\abs{a - b} \leq \abs{a} - \abs{b} \leq \abs{a - b},
      \]
      which is the claim by Theorem~\ref{thm:abs}(c).
    \item $\abs{a - b} = \abs{a + (-b)} \leq \abs{a} + \abs{-b} = \abs{a} + \abs{b}$.\qedhere
  \end{enumerate}
\end{proof}

\begin{remark}
  The trick in (i) is to \emph{add and subtract}: to bound $\abs{a}$ in terms
  of $\abs{a - b}$, write $a = (a - b) + b$ and apply the triangle
  inequality. It will be used constantly.
\end{remark}
